\chapter*{Preface}
\addcontentsline{toc}{chapter}{Preface}
\markboth{Preface}{}
An infinite object can be described by finite rules without being generated by finitely many elements. It can be defined by a short mathematical sentence without admitting computable operations. It can be separated point by point by finite observations without being recoverable from the entire unmarked system of those observations. These distinctions are the subject of this book.

Realization Calculus studies the passage from specified data and admissible operations to an actual mathematical object. Its boundary questions are therefore not external qualifications. They ask which information a realization requires, which operations preserve that information, and which observations discard it. Here these questions are examined in a concrete universe of abelian groups. The constructions are elementary enough to calculate with, yet expressive enough to carry every Turing degree.

The starting object is the direct sum of one cyclic group of each prime order. Its finite stages are cyclic. Its completion is topologically cyclic. The direct sum itself nevertheless needs infinitely many seeds even when every group endomorphism is allowed. Its subgroup lattice then supplies a transparent first model of finite description and effective presentation.

The central construction is more stringent. For every set $A\subseteq\N$, a countable group $G_A$ is placed between that direct sum $D$ and its product completion $P$. All groups have the same torsion subgroup, the same rational quotient, and the same full unmarked finite-quotient system. Nevertheless, their abstract isomorphism types are classified by $A$ modulo finite changes, and an oracle computes a copy precisely when it computes $A$.

\begin{keybox}
\textbf{The central separation.} The compatible finite quotients determine the ambient completion. They do not determine which dense subgroup is the intended original realization. In this family the omitted choice has an exact divisibility signature and an exact computational cost.
\end{keybox}

\section*{How to read the book}
Chapters 1--3 separate the realization contracts and develop the first obstruction examples. Chapters 4--8 give the main construction and its proofs. Chapter 4 also describes the general ambient mechanism: rational subspaces of $P/D$ correspond to pure intermediate groups. Chapters 9--10 extract the reconstruction and grammar consequences. Chapter 11 provides calculations that can be followed by hand and reproduced in Python. Chapter 12 returns to the interpretation in Realization Calculus. The appendices contain implementation guidance, a theorem map, a formalization plan, and a glossary.

Only elementary group theory is needed at the start. The later arguments introduce profinite completion and the relevant computability notions before using them. Readers interested chiefly in the common-completion theorem can begin with Chapters 4--6. Readers interested in its effective meaning should then read Chapters 7--8, since the distinction between a marked and an unmarked group is essential to the decoding proof.

\section*{Evidence and mathematical context}
The main results are supported by the ordinary proofs printed here. The companion programs use exact rational and modular arithmetic to check finite instances and to reject selected false formulas. Those checks do not prove the infinite classification or noncomputability statements. No Lean verification is claimed for this monograph.

The prime-spectrum coding in Chapter 3 has an explicit classical antecedent in Montalb\'an's text \cite{Montalban}. Degrees of isomorphism types, pure subgroups, cocycle extensions, and profinite completion belong to established theories \cite{Calvert,Fuchs,Lupini,Wilkes}. The common-completion family is proved directly rather than attributed to an unverified source. Historical priority for that combined construction has not been established. The contribution presented here is a self-contained theorem development organized around precise realization contracts, not a claim that the classical ingredients originate in this work.

\section*{Conventions that prevent false conclusions}
An unmarked finite-quotient system does not contain the original source group or a distinguished coherent source element. An infinite parameter is not made finite by assigning it a single letter. A computable copy is not the same as an enumerable list of generators and relators. A nonzero divisible quotient does not imply a divisible subgroup. These distinctions appear where they are used, rather than being left to a general disclaimer.
